\chapter{Secure Multiparty Computation}

Every cryptographic chapter so far has had a clean cast: Alice and Bob are allies, Eve is the eavesdropper outside. Multiparty computation breaks that mold. Now the people who want to compute something together \emph{are} the adversaries to one another. They share a goal --- the answer to some function of all their inputs --- but each of them would dearly like to learn everyone else's secrets along the way, and would happily lie or peek if it helped. The astonishing claim of this chapter is that they can get the answer anyway, correctly, while revealing nothing about their private inputs beyond what the answer itself leaks.

That is not a contradiction once you state it carefully. Secure multiparty computation (MPC) lets $n$ parties, holding private inputs $x_1, \dots, x_n$, jointly evaluate an agreed function $f(x_1, \dots, x_n)$ so that everyone learns the output and \emph{nothing else}. No trusted referee, no neutral server holding the plaintext --- just the participants and the math. It is one of the deepest and most counterintuitive results in cryptography, and it powers real systems today: sealed-bid auctions where no auctioneer ever sees the bids, threshold signing where no single machine holds the signing key, and privacy-preserving statistics across hospitals that may not share patient records.

This chapter builds the intuition with the oldest MPC story --- two millionaires comparing their wealth --- then constructs the two great families of protocols: Yao's garbled circuits for two parties, and secret-sharing-based MPC for many. Oblivious transfer turns out to be the small, strange primitive that makes both possible. We keep the numbers tiny so you can trace every wire by hand.

\chapterline

\section{The dream: a trusted party who forgets}
\label{ch34:dream}

Start with what we wish we had. Suppose there were a perfectly trustworthy outside party --- call her the Oracle --- with a magic property: she accepts each participant's secret input, computes the function, announces the result to everyone, and then \emph{instantly forgets} every input she saw. If such an Oracle existed, privacy would be trivial. The millionaires would each whisper their net worth to her; she would announce ``Alice is richer'' or ``Bob is richer'' and nothing more.

MPC's goal is to emulate that Oracle using only messages between the parties themselves, with no Oracle in the room. This is called the \emph{ideal/real paradigm}, and it is the gold standard for defining security here.

\begin{intuitionbox}
The yardstick for ``secure'' is not ``the protocol leaks no information.'' The output itself leaks something: if the auction announces a winning price of \$100, every loser learns their bid was below \$100. The right yardstick is \emph{the protocol leaks nothing more than the ideal Oracle would.} Anything a participant can compute after running the real protocol, they could already have computed from their own input plus the announced output alone. That ``and nothing more'' clause is the entire game.
\end{intuitionbox}

\begin{formalbox}
\textbf{The ideal/real definition (simulation).} Fix a function $f$ and an adversary controlling some subset of corrupt parties. We compare two worlds.

In the \emph{ideal world}, the corrupt parties hand their inputs to the Oracle, who computes $f$ and returns the output. The adversary sees only its own inputs and the output.

In the \emph{real world}, the parties run the actual protocol; the adversary sees every message its corrupt parties send and receive.

The protocol is \emph{secure} if for every real-world adversary there is an ideal-world \emph{simulator} that produces an indistinguishable view using only the ideal-world information (corrupt inputs and the output). If a simulator can fake the entire transcript of received messages knowing only the output, then those messages carried no extra information --- because they could have been manufactured without ever touching the uncorrupted parties' inputs.
\end{formalbox}

We will not prove full simulation security for every protocol below --- the proofs are graduate-course material --- but keep the simulator in mind. Every time we claim ``Bob learns nothing,'' the underlying justification is that Bob's view could have been generated from the output alone.

\subsection{Two adversary models}

Two threat models recur throughout MPC, and they demand different machinery.

\begin{itemize}
  \item \emph{Passive} (also ``curious-but-compliant''): corrupt parties follow the protocol exactly but try to learn extra information from what they legitimately see. Think of a participant who runs the official software but keeps a transcript and studies it later.
  \item \emph{Malicious} (``active''): corrupt parties may deviate arbitrarily --- send wrong messages, abort, lie about inputs --- to break privacy or steer the output.
\end{itemize}

Passive security is much easier to achieve and is the right first target for understanding. Most of this chapter builds passive-secure protocols, then notes what malicious security costs. Real deployments need malicious security, usually bought with extra checks (commitments, zero-knowledge proofs that you followed the rules, authenticated shares).

\begin{connectionbox}
The ``follow the protocol but prove you did'' idea is exactly where zero-knowledge proofs (the chapter on ZK) plug into MPC. A generic compiler due to Goldreich, Micali, and Wigderson upgrades any passive-secure protocol to malicious security by having every party prove in zero knowledge, after each step, that the message they sent is consistent with their committed input and the protocol rules --- without revealing the input. Commitments (the secret-sharing chapter) are the other load-bearing primitive: they pin down inputs so nobody can adaptively change theirs after seeing others' messages.
\end{connectionbox}

\section{The millionaires' problem}
\label{ch34:millionaires}

Andrew Yao posed the founding question in 1982. Two millionaires, Alice and Bob, want to know who is richer without revealing how much either is worth. Alice has wealth $a$, Bob has wealth $b$, both integers in some range $\set{1, \dots, N}$. They want to compute the single bit $[a \ge b]$ and learn nothing else.

\begin{historybox}
Yao's 1982 paper ``Protocols for Secure Computations'' introduced this problem and, in follow-up work, the garbled-circuit technique that solves the general case. The same Andrew Yao won the Turing Award in 2000, partly for this line of work. The millionaires' problem is the ``hello world'' of MPC: small enough to solve by hand, yet it already needs the central trick (oblivious transfer) that the general theory rests on.
\end{historybox}

Here is a clean protocol for the comparison, in the spirit of Yao's original, that already shows the flavor. It assumes Bob has a public-key encryption scheme with the right algebraic structure. We restrict to small ranges so you can follow it; let $a, b \in \set{1, \dots, 10}$.

\begin{enumerate}
  \item Bob picks a random large number $x$, computes $k = \Enc_{pk}(x)$ using \emph{Alice's} public key, and sends Alice the number $k - b$ (the ciphertext interpreted as an integer, minus $b$).
  \item Alice, who holds the secret key, computes $y_u = \Dec_{sk}\!\big((k - b) + u\big)$ for every candidate value $u = 1, 2, \dots, 10$. Exactly one of these, $y_b$, equals $x$, because $(k-b)+b = k$ decrypts to $x$. But Alice does not know which $u$ is $b$.
  \item Alice picks a random prime $p$ of suitable size, reduces every $y_u$ modulo $p$, and constructs a list: for the first $a$ entries ($u = 1, \dots, a$) she sends $y_u \bmod p$ unchanged; for the rest ($u = a+1, \dots, 10$) she sends $(y_u + 1) \bmod p$. She sends Bob the whole list of $10$ numbers, plus $p$.
  \item Bob looks at position $b$ in the list and compares it against $x \bmod p$. If the $b$-th entry equals $x \bmod p$, then position $b$ fell in the ``unchanged'' region, meaning $b \le a$, so $a \ge b$. Otherwise $a < b$.
\end{enumerate}

\begin{walkbox}
Why does Bob learn only the one bit? Every number Alice sends is either $y_u \bmod p$ or $(y_u + 1) \bmod p$, and the $y_u$ for $u \ne b$ are decryptions of values unrelated to $x$ --- effectively random mod $p$ from Bob's view, since he does not hold the secret key. So the list looks like noise to him except at his own position $b$, where he can recognize $x \bmod p$. He reads off whether his entry was bumped by $1$ or not. That single comparison is the bit $[b \le a]$. He cannot tell $a$'s exact value, only which side of $b$ it lies on.

Why does Alice learn nothing? She never sees $b$. She decrypts $10$ candidates and one happens to be $x$, but she cannot tell which, because $x$ is random and Bob's offset $k - b$ hides $b$ inside the ciphertext arithmetic.
\end{walkbox}

\begin{trapbox}
This particular protocol is delicate: it relies on the encryption being well-behaved under the integer subtraction $k - b$, which is not true for a generic modern encryption scheme. Yao's original used a number-theoretic scheme where it worked. Do not implement comparisons this way in production; the point here is to \emph{feel} how two mutually distrustful parties can squeeze out one bit. The general, robust solution is garbled circuits, which we build next, and which compute the comparison --- or any function --- without such scheme-specific tricks.
\end{trapbox}

\section{Oblivious transfer: the atom of MPC}
\label{ch34:ot}

Before circuits, we need the primitive everything is made of. \emph{Oblivious transfer} (OT) sounds mundane and turns out to be the universal building block: with OT you can build \emph{any} secure two-party computation.

The basic flavor is \emph{1-out-of-2 OT}. The sender holds two messages $m_0$ and $m_1$. The receiver holds a choice bit $c \in \set{0,1}$. After the protocol:

\begin{itemize}
  \item the receiver learns $m_c$ (the message they chose) and nothing about $m_{1-c}$;
  \item the sender learns nothing --- in particular, not which message the receiver picked.
\end{itemize}

\begin{intuitionbox}
Picture two sealed boxes on a table, $m_0$ and $m_1$. The receiver gets to walk away with exactly one box, of their choosing, and cannot peek inside the other. The sender, watching, cannot tell which box was taken. It feels paradoxical: how can the sender hand over one of two secrets without knowing which, while keeping the other hidden? Public-key math makes it possible.
\end{intuitionbox}

\subsection{A concrete OT from Diffie-Hellman}

Here is a clean OT protocol (a simplified Bellare-Micali / Naor-Pinkas style construction) over a cyclic group $\G = \langle g \rangle$ of prime order $q$, where the Diffie-Hellman problem is hard. Recall from the discrete-log chapters that knowing $g^x$ and $g^y$ does not let you compute $g^{xy}$.

\textbf{Setup.} Public group $\G$, generator $g$, and a fixed group element $C$ whose discrete log nobody knows (a ``nothing-up-my-sleeve'' constant).

\textbf{Protocol.}
\begin{enumerate}
  \item Receiver (choice bit $c$) picks random $k \in \set{0,\dots,q-1}$. He sets $PK_c = g^k$ and $PK_{1-c} = C / g^k = C \cdot g^{-k}$. He sends $PK_0$ to the sender.
  \item The sender computes $PK_1 = C / PK_0$ herself. Now she has both ``public keys'' $PK_0$ and $PK_1$, with the property $PK_0 \cdot PK_1 = C$. She does \emph{not} know which one the receiver can decrypt.
  \item The sender picks random $r_0, r_1$ and encrypts: she sends
  \[
    \big(g^{r_0},\ m_0 \xor H(PK_0^{\,r_0})\big) \quad\text{and}\quad \big(g^{r_1},\ m_1 \xor H(PK_1^{\,r_1})\big),
  \]
  where $H$ hashes a group element to a bit string.
  \item The receiver knows $k$ such that $PK_c = g^k$. For his chosen branch he computes the masking key $H\big((g^{r_c})^{k}\big) = H(PK_c^{\,r_c})$ and unmasks $m_c$.
\end{enumerate}

\begin{whybox}
\textbf{Why the receiver gets exactly one message.} The receiver knows the discrete log $k$ of $PK_c$ only. To unmask $m_{1-c}$ he would need $PK_{1-c}^{\,r_{1-c}} = (g^{r_{1-c}})^{\log PK_{1-c}}$, i.e.\ he would need the discrete log of $PK_{1-c} = C \cdot g^{-k}$. That requires the discrete log of $C$, which nobody knows. So the other message stays a sealed box --- this is exactly the Diffie-Hellman hardness assumption doing the work.

\textbf{Why the sender stays oblivious.} The only thing the sender sees from the receiver is $PK_0$. Since $PK_0 = g^k$ for a uniformly random $k$, it is a uniformly random group element regardless of whether $c = 0$ (so $PK_0 = g^k$) or $c = 1$ (so $PK_0 = C \cdot g^{-k}$, also uniform because $g^{-k}$ is uniform). The sender's view is identical in both cases. She cannot learn $c$.
\end{whybox}

\begin{labbox}
\textbf{Trace it with tiny numbers.} Work in $\Zns{23}$, restricted to the subgroup of order $q = 11$ generated by $g = 2$ (since $2^{11} \equiv 1 \pmod{23}$). Let the public constant be $C = 2^6 \bmod 23 = 18$ (pretend nobody knows the exponent $6$). Suppose the receiver's choice bit is $c = 1$ and he picks $k = 4$.

He sets $PK_1 = g^k = 2^4 = 16$ and $PK_0 = C / g^k = 18 \cdot 16^{-1} \bmod 23$. Now $16^{-1} \bmod 23$: since $16 \cdot 13 = 208 = 9\cdot 23 + 1$, the inverse is $13$. So $PK_0 = 18 \cdot 13 = 234 \equiv 4 \pmod{23}$. He sends $PK_0 = 4$.

The sender computes $PK_1 = C / PK_0 = 18 \cdot 4^{-1}$. Here $4^{-1} \bmod 23 = 6$ (since $4\cdot 6 = 24 \equiv 1$), so $PK_1 = 18 \cdot 6 = 108 \equiv 16 \pmod{23}$ --- matching the receiver's $PK_1 = 16$. Check: $PK_0 \cdot PK_1 = 4 \cdot 16 = 64 \equiv 18 = C$. Good. The sender now encrypts $m_0$ under $PK_0$ and $m_1$ under $PK_1$ with fresh $r_0, r_1$; only the $m_1$ branch can be opened, because the receiver knows $k=4$ is the log of $PK_1$ but has no clue about the log of $PK_0 = 4$.
\end{labbox}

\begin{focusbox}
OT is to MPC what AND/NOT gates are to circuits: a complete basis. Kilian proved in 1988 that OT alone is enough to build secure computation of \emph{any} function. So once you have a working OT --- from Diffie-Hellman, from RSA, from a quantum channel, whatever --- you have, in principle, secure computation of everything. The rest of two-party MPC is about doing it efficiently.
\end{focusbox}

\subsection{OT extension: making it cheap}

Public-key operations are expensive, and a circuit may need millions of OTs (one per input wire of one party). The breakthrough that made MPC practical is \emph{OT extension} (Ishai-Kilian-Nissim-Petrank, 2003): perform a small fixed number of ``real'' public-key OTs (say $128$), then use cheap symmetric cryptography (hashing) to bootstrap those into millions of OTs. It is the OT analogue of using public-key crypto only to set up a symmetric session key, then encrypting the bulk traffic with AES. Without OT extension, garbled-circuit MPC would be a curiosity; with it, you can evaluate circuits with billions of gates.

\section{Yao's garbled circuits}
\label{ch34:garbled}

Now the centerpiece of two-party MPC. The idea: take the function $f$ and express it as a Boolean circuit of AND, XOR, and NOT gates. One party (the \emph{garbler}, say Alice) ``encrypts'' the entire circuit --- she garbles it --- so that the other party (the \emph{evaluator}, Bob) can run it on hidden inputs and read off only the final output, learning nothing about the intermediate wire values or Alice's input.

\subsection{Garbling one wire, one gate}

Every wire in the circuit carries a bit, $0$ or $1$. In a garbled circuit, Alice replaces each bit value on each wire with a random \emph{label} --- a fresh $128$-bit key. For wire $w$ she picks two labels: $W^0$ (meaning ``this wire is $0$'') and $W^1$ (meaning ``this wire is $1$''). The labels look random; seeing one tells you nothing about whether it encodes $0$ or $1$.

Now consider a single AND gate with input wires $a, b$ and output wire $z$. The truth table is:

\[
\begin{array}{cc|c}
a & b & z = a \land b \\\hline
0 & 0 & 0 \\
0 & 1 & 0 \\
1 & 0 & 0 \\
1 & 1 & 1
\end{array}
\]

Alice has labels $A^0, A^1$ for wire $a$; $B^0, B^1$ for wire $b$; $Z^0, Z^1$ for wire $z$. She builds a \emph{garbled table}: for each of the four input combinations, she encrypts the correct output label under the two corresponding input labels as keys. Using a double-key encryption $\Enc_{K_1, K_2}(\cdot)$ (think: encrypt under $K_1$, then under $K_2$):

\[
\begin{aligned}
&\Enc_{A^0, B^0}(Z^0) \\
&\Enc_{A^0, B^1}(Z^0) \\
&\Enc_{A^1, B^0}(Z^0) \\
&\Enc_{A^1, B^1}(Z^1)
\end{aligned}
\]

Alice randomly \emph{shuffles} these four rows (this is the ``garbling'' --- the row order no longer leaks the input bits) and sends the scrambled table to Bob.

\begin{whybox}
\textbf{Why Bob can evaluate but not snoop.} Suppose Bob holds exactly one input label for each input wire --- say $A^1$ (so wire $a = 1$, but Bob does not know that) and $B^0$. He tries to decrypt each of the four rows with the key pair $(A^1, B^0)$. Three rows were encrypted under different key pairs, so they decrypt to garbage (a well-designed double encryption lets Bob detect this --- for example, by a block of zero padding that only appears on correct decryption). Exactly one row, the $\Enc_{A^1,B^0}(Z^0)$ row, decrypts cleanly, and Bob recovers the output label $Z^0$.

Bob now holds $Z^0$, the label for ``output is $0$'' --- and indeed $1 \land 0 = 0$. But $Z^0$ is just a random $128$-bit string to him; he does not know it means $0$. He carries it forward as the input label to the next gate. He never learns the actual bit on any internal wire. He only ever holds one label per wire, never both, so he can never decrypt more than one row and never reconstruct the table.
\end{whybox}

\begin{figure}[H]
\centering
\begin{tikzpicture}[font=\small]
  \node (a) [draw, rounded corners, fill=aliceblue!12, minimum width=22mm] {wire $a$: $A^0, A^1$};
  \node (b) [draw, rounded corners, fill=aliceblue!12, minimum width=22mm, below=10mm of a] {wire $b$: $B^0, B^1$};
  \node (g) [draw, fill=structurepurple!15, minimum width=22mm, minimum height=14mm, right=22mm of $(a)!0.5!(b)$] {AND};
  \draw[-Latex] (a.east) -- (g.west |- a);
  \draw[-Latex] (b.east) -- (g.west |- b);
  \node (z) [draw, rounded corners, fill=targetgreen!15, minimum width=22mm, right=22mm of g] {wire $z$: $Z^0, Z^1$};
  \draw[-Latex] (g.east) -- (z.west);
  \node (t) [draw, fill=codebg, align=left, below=14mm of g, font=\ttfamily\footnotesize] {
    Enc(A1,B1 -> Z1)\\
    Enc(A0,B0 -> Z0)\\
    Enc(A1,B0 -> Z0)\\
    Enc(A0,B1 -> Z0)};
  \node[below=1mm of t, font=\footnotesize\itshape] {garbled \& shuffled table sent to Bob};
  \draw[-Latex, gray] (g.south) -- (t.north);
\end{tikzpicture}
\caption{One garbled AND gate. Alice replaces each wire bit with a random label, encrypts each output label under its two input labels, and shuffles the rows. Bob, holding one label per input wire, decrypts exactly one row and learns one output label --- never the underlying bits.}
\label{ch34:fig-garbled-gate}
\end{figure}

\subsection{Wiring up the whole circuit}

A circuit is many gates chained together. Alice garbles every gate the same way, so the output label of one gate becomes an input label of the next. The remaining questions are: how does Bob get the input labels to start, and how does he read the final answer?

\begin{enumerate}
  \item \textbf{Alice's input.} For each of Alice's input wires, she knows her bit, so she just hands Bob the corresponding label directly. If Alice's input bit on wire $a$ is $1$, she sends $A^1$. Bob cannot tell it means $1$.
  \item \textbf{Bob's input --- via OT.} Bob must get the label matching \emph{his} input bit on each of his input wires, but Alice must not learn which bit Bob has, and Bob must not get both labels. This is exactly 1-out-of-2 oblivious transfer. For each of Bob's input wires, Alice plays OT sender with messages $(B^0, B^1)$; Bob plays receiver with choice bit equal to his input. He gets $B^{c}$ and nothing else, and Alice learns nothing about $c$.
  \item \textbf{Evaluation.} Bob propagates labels gate by gate through the garbled tables, decrypting one row per gate, until he holds the labels on the output wires.
  \item \textbf{Output decoding.} Alice provides a small decoding table mapping the two output-wire labels back to bits: ``$Z^0_{\text{out}} \mapsto 0$, $Z^1_{\text{out}} \mapsto 1$.'' Bob looks up his final label and reads the answer. (If only Bob should learn the output, Alice sends him this table privately; if both should, Bob announces the labels and Alice decodes too.)
\end{enumerate}

\begin{figure}[H]
\centering
\begin{tikzpicture}[font=\small]
  \node[draw, fill=aliceblue!12, minimum width=20mm] (alice) {Alice (garbler)};
  \node[draw, fill=bobgreen!12, minimum width=20mm, right=58mm of alice] (bob) {Bob (evaluator)};
  \draw[-Latex] ([yshift=-7mm]alice.south) -- node[above, font=\footnotesize] {garbled circuit + Alice's labels} ([yshift=-7mm]bob.south);
  \draw[-Latex] ([yshift=-16mm]alice.south) -- node[above, font=\footnotesize] {OT: Bob's input labels} node[below, font=\footnotesize\itshape] {(Alice cannot see Bob's bits)} ([yshift=-16mm]bob.south);
  \draw[-Latex] ([yshift=-25mm]alice.south) -- node[above, font=\footnotesize] {output decoding table} ([yshift=-25mm]bob.south);
  \draw[Latex-, gray] ([yshift=-32mm]alice.south) -- node[below, font=\footnotesize] {(optional) output labels back to Alice} ([yshift=-32mm]bob.south);
  \draw[dashed, gray] (alice.south) -- ([yshift=-34mm]alice.south);
  \draw[dashed, gray] (bob.south) -- ([yshift=-34mm]bob.south);
\end{tikzpicture}
\caption{The message flow of Yao's two-party protocol. Alice garbles and sends the circuit; Bob obtains his own input labels through oblivious transfer; Bob evaluates locally and decodes the output.}
\label{ch34:fig-yao-flow}
\end{figure}

\begin{whybox}
\textbf{Why the whole thing is secure (passive model).} Bob's entire view is: a pile of ciphertexts (the garbled tables), one label per wire, and a decoding table. A simulator who knows only the final output can fabricate an indistinguishable view --- it generates random labels, builds fake garbled tables that decrypt consistently along the single path Bob walks, and rigs the decoding table to yield the known output. Because Bob only ever sees one label per wire and the tables are encrypted, he cannot distinguish the real garbling from the simulated one. Hence he learns nothing beyond the output. Symmetrically, Alice sees only the OT transcript, from which (by OT's sender-obliviousness) she learns nothing about Bob's input.
\end{whybox}

\subsection{The optimizations that make it real}

Plain garbling sends four ciphertexts per gate. Decades of refinement shrank this dramatically, and these tricks are why garbled circuits run at line rate today.

\begin{itemize}
  \item \emph{Point-and-permute}: append a random ``select bit'' to each label so Bob knows directly which single row to decrypt, instead of trial-decrypting all four. One decryption per gate.
  \item \emph{Free XOR} (Kolesnikov-Schneider, 2008): choose all labels so that $W^1 = W^0 \xor \Delta$ for a single global secret offset $\Delta$. Then an XOR gate needs \emph{no} garbled table and \emph{no} cryptographic operation at all --- the output label is just the XOR of the input labels. XOR gates become free, which matters because many useful circuits (addition, comparison) are XOR-heavy.
  \item \emph{Half-gates} (Zahur-Rosulek-Evans, 2015): brings each AND gate down to two ciphertexts, proven optimal for a broad class of techniques.
\end{itemize}

\begin{implementbox}
A stripped-down garbler for a single AND gate, to make the labels and tables concrete. This is \textbf{toy code for learning only} --- it uses a plain SHA-256 mask, no authentication beyond a toy membership check, no point-and-permute, and leaks timing. Never use it for anything real; production garbling uses AES with hardware acceleration and the optimizations above.

\begin{lstlisting}[language=Python]
import os, hashlib, random

def label():                      # a random 16-byte wire label
    return os.urandom(16)

def kdf(k1, k2):                  # derive a mask from two key labels
    return hashlib.sha256(k1 + k2).digest()[:16]

def xor(a, b):
    return bytes(x ^ y for x, y in zip(a, b))

# two labels per wire: index 0 means bit 0, index 1 means bit 1
A = [label(), label()]
B = [label(), label()]
Z = [label(), label()]

# AND truth table: z = a AND b
table = []
for a in (0, 1):
    for b in (0, 1):
        z = a & b
        ct = xor(Z[z], kdf(A[a], B[b]))   # "encrypt" output label
        table.append(ct)
random.shuffle(table)                      # garble: hide row order

# --- Bob evaluates with one label per wire, say a=1, b=0 ---
my_A, my_B = A[1], B[0]
recovered = None
for ct in table:
    cand = xor(ct, kdf(my_A, my_B))
    if cand in (Z[0], Z[1]):               # toy "did it decrypt?" check
        recovered = cand
print("output label == Z0 (bit 0)?", recovered == Z[0])  # True: 1 AND 0 = 0
\end{lstlisting}
Real systems replace the membership test with an authenticated decryption (a row succeeds only with valid padding or a MAC), so Bob cannot brute-force other rows.
\end{implementbox}

\section{Secret-sharing-based MPC for many parties}
\label{ch34:ss-mpc}

Garbled circuits are a two-party masterpiece, but they do not scale gracefully to $n$ parties. For many participants, the dominant approach is built on \emph{secret sharing} --- the Shamir scheme from the commitments chapter. The idea is direct: split every secret into shares, then compute on the shares.

\subsection{Recap: Shamir shares are linear}

Recall Shamir $t$-out-of-$n$ sharing. To share a secret $s \in \F_p$, the dealer picks a random degree-$(t-1)$ polynomial $f$ with $f(0) = s$, and gives party $i$ the share $f(i)$. Any $t$ parties can interpolate to recover $s$; any $t-1$ learn nothing, because $t-1$ points are consistent with every possible secret.

The magic property for MPC: \emph{secret sharing is linear}. If $f$ shares $a$ and $g$ shares $b$, then $f + g$ is a valid polynomial sharing $a + b$ --- and each party can compute its share of the sum locally, just by adding its two shares. No communication needed.

\begin{intuitionbox}
Think of each party as holding a single coordinate of a high-dimensional secret. Adding two secrets means adding coordinates pointwise; nobody needs to see anyone else's coordinate to do that. Scaling by a public constant is just as local: multiply your share by the constant. This is why \emph{linear} functions of secrets are trivial in this model. The whole difficulty of MPC concentrates in one place: multiplication.
\end{intuitionbox}

\subsection{Addition is free, multiplication is the hard part}

Let each party $P_i$ hold a Shamir share $[x] = (x_1, \dots, x_n)$ of secret $x$ and $[y]$ of $y$, where $[x]$ means ``the vector of shares of $x$.''

\textbf{Addition.} To get shares of $x + y$, each party computes $x_i + y_i$ locally. Done. Shares of $x + c$ for a public constant $c$: by convention one designated party adds $c$ to its share. Shares of $c \cdot x$: each party multiplies its share by $c$.

\textbf{Multiplication.} Each party can multiply locally: $x_i \cdot y_i$. But there is a problem. If $f$ has degree $t-1$ and $g$ has degree $t-1$, then their product $f \cdot g$ has degree $2(t-1)$. The local products $x_i \cdot y_i$ are points on this degree-$2(t-1)$ polynomial, which does share $xy$ at $0$ --- but you now need $2t-1$ points to interpolate it, and after a few multiplications the degree blows past $n$ and the secret becomes unrecoverable. Worse, $f \cdot g$ is not a uniformly random polynomial of that degree (its coefficients are correlated), so the naive product can leak. Multiplication needs a \emph{degree-reduction} step with communication.

\begin{figure}[H]
\centering
\begin{tikzpicture}[font=\small]
  \draw[->] (0,0) -- (6.4,0) node[right] {party index $i$};
  \draw[->] (0,0) -- (0,3.2) node[above] {share value};
  \draw[aliceblue, thick] (0.3,0.6) -- (5.6,2.6) node[right, black] {$[x]$};
  \draw[bobgreen, thick] (0.3,2.4) -- (5.6,0.8) node[right, black] {$[y]$};
  \draw[trapred, thick, domain=0.3:5.6, samples=40, smooth] plot (\x, {0.13*(\x-3)*(\x-3)+0.5}) node[right, black] {$[xy]$};
  \node[trapred, font=\footnotesize\itshape] at (3,3.0) {degree doubled --- must reduce};
\end{tikzpicture}
\caption{Why multiplication is hard. The shares $[x],[y]$ lie on degree-$(t{-}1)$ polynomials, but the local products lie on a polynomial of doubled degree (and skewed distribution). A re-sharing round both reduces the degree back to $t-1$ and re-randomizes it.}
\label{ch34:fig-degree}
\end{figure}

\subsection{The BGW multiplication step}

The classic fix (Ben-Or, Goldwasser, Wigderson, 1988) handles the multiplication with one round of communication:

\begin{enumerate}
  \item Each party $P_i$ computes the local product $d_i = x_i \cdot y_i$. These $d_i$ lie on a degree-$2(t-1)$ polynomial $h$ with $h(0) = xy$.
  \item Each $P_i$ \emph{re-shares} its $d_i$: it acts as a dealer and Shamir-shares $d_i$ with a fresh degree-$(t-1)$ polynomial, sending sub-shares to everyone.
  \item Now everyone holds shares of every $d_i$. They take a fixed public linear combination of these (the Lagrange coefficients that would reconstruct $h(0)$) --- but applied \emph{on the shares}, which is a local linear operation. The result is a fresh degree-$(t-1)$ sharing of $\sum_i \lambda_i d_i = h(0) = xy$.
\end{enumerate}

The interpolation that would normally combine the $d_i$ into $xy$ is a linear map, and we already know linear maps on shares are free. By pushing it through the re-shared low-degree polynomials, we get $xy$ shared at degree $t-1$ again, re-randomized. This needs $t < n/2$ (a majority of uncorrupted parties) for the degree-$2(t-1)$ polynomial to be recoverable.

\begin{focusbox}
The pattern is universal: \emph{linear operations are local and free; multiplication costs one interaction.} So the ``size'' of an MPC is governed by the number of multiplication gates in the circuit (its \emph{multiplicative depth} and count), exactly as garbled circuits are governed by AND gates and free on XOR. Designing MPC-friendly functions means minimizing multiplications --- which is why MPC and homomorphic encryption (the previous chapter) share so much circuit-optimization lore.
\end{focusbox}

\subsection{A worked multiplication with tiny numbers}

Work over $\F_{11}$ (arithmetic mod $11$) with $n = 3$ parties and threshold $t = 2$, so polynomials are degree $1$ (lines). We multiply $x = 3$ by $y = 4$; the true answer is $12 \equiv 1 \pmod{11}$.

\textbf{Share $x = 3$.} Dealer picks the line $f(u) = 3 + 2u$. Shares: $f(1) = 5$, $f(2) = 7$, $f(3) = 9$. So $x_1, x_2, x_3 = 5, 7, 9$.

\textbf{Share $y = 4$.} Dealer picks the line $g(u) = 4 + 1u$. Shares: $g(1) = 5$, $g(2) = 6$, $g(3) = 7$. So $y_1, y_2, y_3 = 5, 6, 7$.

\textbf{Local products.} $d_1 = 5\cdot 5 = 25 \equiv 3$, $d_2 = 7\cdot 6 = 42 \equiv 9$, $d_3 = 9\cdot 7 = 63 \equiv 8 \pmod{11}$. These three points $(1,3),(2,9),(3,8)$ lie on a degree-$2$ parabola $h$ with $h(0) = xy$.

\textbf{Check the answer is hidden in $h(0)$.} Interpolate $h$ through the three points using Lagrange at $u = 0$. The Lagrange coefficients at $0$ for nodes $1,2,3$ are
\[
\lambda_1 = \tfrac{(0-2)(0-3)}{(1-2)(1-3)} = \tfrac{6}{2} = 3,\quad
\lambda_2 = \tfrac{(0-1)(0-3)}{(2-1)(2-3)} = \tfrac{3}{-1} = -3,\quad
\lambda_3 = \tfrac{(0-1)(0-2)}{(3-1)(3-2)} = \tfrac{2}{2} = 1.
\]
Then $h(0) = \lambda_1 d_1 + \lambda_2 d_2 + \lambda_3 d_3 = 3\cdot 3 + (-3)\cdot 9 + 1\cdot 8 = 9 - 27 + 8 = -10 \equiv 1 \pmod{11}$. That is exactly $xy = 1$. The product is correctly encoded --- it just sits on a degree-$2$ polynomial. The BGW re-sharing step would now push this same linear combination through fresh degree-$1$ sharings of each $d_i$, leaving everyone with a degree-$1$ sharing of $1$, ready for the next gate.

\subsection{Beaver triples: precompute the hard part}

A beautiful refinement (Donald Beaver, 1991) moves the expensive multiplication \emph{offline}. In a preprocessing phase --- before inputs are even known --- the parties generate shared random \emph{triples} $([a], [b], [c])$ with $c = ab$, where $a, b$ are random and unknown to anyone. Then, in the online phase, multiplying two shared secrets $[x]\cdot[y]$ costs only additions and two public reconstructions:

\begin{enumerate}
  \item Parties locally compute and reveal $[d] = [x] - [a]$ and $[e] = [y] - [b]$. Since $a, b$ are random and secret, $d = x - a$ and $e = y - b$ leak nothing about $x, y$ (they are one-time-pad-masked).
  \item Each party locally sets its share of $[xy]$ to
  \[
    [xy] = [c] + d\,[b] + e\,[a] + d\,e.
  \]
\end{enumerate}

\begin{whybox}
\textbf{Why the triple formula gives $xy$.} Substitute $x = d + a$ and $y = e + b$:
\[
xy = (d+a)(e+b) = de + d b + e a + ab = de + d\,b + e\,a + c.
\]
Every term on the right is either a public value ($d, e$) times a share, or the precomputed $[c]$ --- all linear operations on shares, hence local and free. The only communication in the online phase is revealing $d$ and $e$, two field elements per multiplication. The secret-protecting masking is exactly a one-time pad: $d = x - a$ hides $x$ because $a$ is uniform and used once.
\end{whybox}

\begin{focusbox}
This offline/online split is the architecture of modern practical MPC frameworks (SPDZ and its descendants). All the cryptographically heavy lifting --- generating authenticated multiplication triples, often via homomorphic encryption or OT extension --- happens in advance, in bulk, when there is time. When the real inputs arrive, the online phase is almost pure information-theoretic arithmetic and runs very fast. SPDZ also adds message authentication codes on every share so that a \emph{malicious} party who tampers with a value is caught with overwhelming probability, lifting security from the passive model to the malicious one.
\end{focusbox}

\section{An attack: when ``private'' inputs aren't enough}
\label{ch34:attack}

MPC's guarantee is precise and easy to over-read. It protects the inputs \emph{during} computation. It says nothing about whether the agreed function and output are safe to reveal. This gap is where real deployments get hurt.

\begin{attackbox}
\textbf{The function-leakage attack.} Suppose three hospitals run an MPC to compute the average age of patients with a rare disease, to publish a statistic. The protocol is flawless: no hospital sees another's records. But suppose hospital $A$ controls (or colludes with) two of the three parties and knows its own contribution exactly. From the published average and its own data, $A$ can algebraically solve for the combined contribution of the others. If only one other patient exists, $A$ learns that individual's exact age --- a total privacy breach, with the MPC working perfectly. The leak came from the \emph{function}, not the protocol.

\textbf{Even nastier: the chosen-input attack.} A malicious party can pick a \emph{crafted} input to make the output reveal a target's secret. To learn whether Bob's salary equals \$87{,}000, a malicious participant in a ``compute the maximum salary'' MPC submits the input \$86{,}999. If the announced max is \$87{,}000, and the attacker knows everyone else is below \$86{,}999, the output pinpoints Bob. Again the cryptography did its job; the function composed with adversarial inputs did the leaking.
\end{attackbox}

\begin{trapbox}
Two defenses, neither cryptographic in the usual sense:
\begin{itemize}
  \item \emph{Function vetting / output privacy}: only compute functions whose output is safe to publish. This is where MPC meets \emph{differential privacy} --- add calibrated noise inside the MPC so the output cannot single out an individual. MPC hides inputs; differential privacy bounds what the output reveals. Serious privacy systems need both.
  \item \emph{Input validation and collusion bounds}: enforce that inputs are well-formed (via zero-knowledge range proofs) and design for a realistic corruption threshold. A protocol that assumes most parties are uncorrupted is worthless if one entity secretly controls a majority of the ``independent'' parties.
\end{itemize}
The lesson: MPC is a tool for input privacy, not a magic wand for ``privacy.'' Always ask separately whether the \emph{output} is safe.
\end{trapbox}

\section{Where MPC is used in the real world}
\label{ch34:applications}

\begin{applicationbox}
\textbf{Sealed-bid auctions.} The first large-scale real MPC deployment was Denmark's 2008 sugar-beet auction (Bogetoft et al.). Thousands of farmers and the sole buyer needed a market-clearing price, but neither side trusted the other --- or any third party --- to see the bids, which revealed each farm's economics. Three servers ran a secret-sharing MPC; no machine ever held a plaintext bid, yet the clearing price came out correctly. This is the canonical ``the auctioneer never sees the bids'' application.

\textbf{Threshold signing and key management.} A signing key split across $n$ machines via secret sharing, where any $t$ can jointly produce a signature but no smaller set can, and no machine ever reconstructs the full key. Threshold ECDSA and threshold Schnorr/EdDSA protocols are MPC for the specific function ``sign this message.'' Cryptocurrency custodians, certificate authorities, and HSM clusters use this so that compromising one server (or one insider) cannot steal the key or forge signatures. The elliptic-curve machinery from the ECC chapter is exactly what these protocols compute over, securely shared.

\textbf{Privacy-preserving analytics.} Google's ``Private Join and Compute'' and similar private-aggregation systems let parties compute aggregate statistics (ad-conversion counts, telemetry sums) across datasets without sharing the underlying records. Apple and Google's exposure-notification and private-set-intersection features use OT/MPC ideas so one party learns the intersection size or a sum without learning the other's full set.

\textbf{Private set intersection (PSI).} A specialized, hyper-optimized MPC: two parties learn only the elements their sets share. ``Contact discovery'' in messengers (find which of my contacts use the app) and ``did this leaked password appear in your breach database'' lookups are PSI in production, built largely from OT extension.
\end{applicationbox}

\begin{connectionbox}
MPC, homomorphic encryption (previous chapter), and zero-knowledge proofs form a triad of ``computing on secrets.'' Homomorphic encryption lets \emph{one} party compute on another's encrypted data --- great when computation is outsourced to a single untrusted server. MPC distributes the computation across the data owners themselves --- great when there is no single trusted server and inputs come from many parties. ZK proves a statement about a secret without revealing it --- and is the glue that makes MPC robust against malicious parties. Real systems mix them: SPDZ uses homomorphic encryption to generate Beaver triples; threshold signatures use ZK proofs to keep parties in line.
\end{connectionbox}

\section{What it costs, and when to reach for it}
\label{ch34:costs}

MPC is not free. Garbled circuits send ciphertexts proportional to the number of AND gates and require a full circuit representation of $f$ --- expensive for functions with deep arithmetic. Secret-sharing MPC trades that for rounds of communication proportional to the multiplicative depth, so high-latency networks hurt. Both are orders of magnitude slower than computing $f$ in the clear, though OT extension, free-XOR, half-gates, and the SPDZ offline/online split have pushed many useful computations into the practical range (seconds to minutes for serious circuits, microseconds per gate).

\begin{marginnotebox}
Rough mental model for choosing a technique:
\begin{itemize}
  \item Two parties, function expressible as a Boolean circuit, want low round count: \emph{Yao garbled circuits} (constant rounds).
  \item Many parties, arithmetic-heavy function, a majority of uncorrupted parties, can tolerate rounds: \emph{secret-sharing MPC} (BGW/SPDZ).
  \item Need malicious security: add MACs (SPDZ) or ZK compilation; expect a constant-factor to order-of-magnitude slowdown.
  \item One untrusted compute server, not many data owners: consider \emph{homomorphic encryption} instead (previous chapter).
\end{itemize}
\end{marginnotebox}

The headline to carry away: any function that can be computed can be computed \emph{securely} among mutually distrustful parties, with privacy of inputs guaranteed by oblivious transfer (for circuits) or by the linearity of secret sharing (for arithmetic). The remaining engineering is about making that theoretical possibility fast enough to use --- and about the separate, non-cryptographic question of whether the output you chose to reveal was safe to reveal at all.

\section*{Exercises}
\addcontentsline{toc}{section}{Exercises}

\begin{exercise}
\textbf{OT by hand.} Use the Diffie-Hellman OT of Section~\ref{ch34:ot} in the order-$11$ subgroup of $\Zns{23}$ with $g = 2$ and public constant $C = 18$. Take the receiver's choice bit $c = 0$ and $k = 5$. Compute $PK_0$ and $PK_1$, verify $PK_0 \cdot PK_1 \equiv C \pmod{23}$, and state which of the two messages the receiver can recover and why he cannot recover the other.
\end{exercise}

\begin{exercise}
\textbf{Garble an OR gate.} Following Section~\ref{ch34:garbled}, write out the four garbled-table rows for an OR gate $z = a \lor b$ in terms of the wire labels $A^0, A^1, B^0, B^1, Z^0, Z^1$. Then, given that Bob holds $A^0$ and $B^1$, identify which row he decrypts and which output label he obtains. What bit does that label encode?
\end{exercise}

\begin{exercise}
\textbf{Free XOR.} Suppose all labels satisfy $W^1 = W^0 \xor \Delta$ for a global secret $\Delta$. Show that if the output label of an XOR gate is defined as $Z^0 = A^0 \xor B^0$ (and $Z^1 = Z^0 \xor \Delta$), then for \emph{any} input bits $a, b$ the evaluator who holds $A^a$ and $B^b$ can compute the correct output label $A^a \xor B^b$ with no garbled table. Verify all four cases.
\end{exercise}

\begin{exercise}
\textbf{Local addition is correct.} Over $\F_{11}$ with $t = 2$, $n = 3$, share $x = 6$ with the line $f(u) = 6 + 3u$ and $y = 2$ with $g(u) = 2 + 5u$. Write each party's two shares, have each party add them locally, and interpolate the three sums back to $u = 0$. Confirm you recover $x + y = 8$ and that the parties needed no communication.
\end{exercise}

\begin{exercise}
\textbf{Multiplication degree blow-up.} Using the same shares of $x = 6$ and $y = 2$ from the previous exercise, compute the three local products $d_i = x_i y_i \bmod 11$. Interpolate the parabola through $(1,d_1),(2,d_2),(3,d_3)$ and evaluate at $0$ using the Lagrange coefficients $3, -3, 1$. Confirm you get $xy = 12 \equiv 1 \pmod{11}$, and explain why the parties cannot simply use the $d_i$ directly as shares for the next gate.
\end{exercise}

\begin{exercise}
\textbf{Beaver triple arithmetic.} Over $\F_{13}$, a precomputed triple has $a = 5$, $b = 4$, $c = ab = 20 \equiv 7$. Parties hold shares of $x = 9$ and $y = 6$. Compute the revealed values $d = x - a$ and $e = y - b$ modulo $13$, then evaluate $c + db + ea + de$ modulo $13$ and check it equals $xy \bmod 13$. (You may work with the secrets directly to verify the formula; the point is the algebra.)
\end{exercise}

\begin{exercise}
\textbf{Break a careless deployment.} Five companies run an MPC to publish the \emph{sum} of their R\&D budgets. Four of them quietly collude. Explain precisely what they learn about the fifth company's budget, and why the MPC protocol itself is not at fault. Propose one mitigation that does not abandon the MPC.
\end{exercise}

\begin{exercise}
\textbf{Chosen-input attack.} A salary-transparency MPC computes the \emph{maximum} salary among participants and announces it. You are malicious and want to learn whether a specific colleague earns exactly \$120{,}000, and you believe everyone else earns less than \$119{,}000. Describe the input you submit and how you read off the answer from the announced maximum. What property of the chosen \emph{function} enabled the attack?
\end{exercise}

\begin{exercise}
\textbf{Why OT needs hardness.} Argue that 1-out-of-2 oblivious transfer cannot exist with information-theoretic security against \emph{both} parties simultaneously in the plain two-party model (no setup). Hint: consider what it would mean for the receiver to be unable to learn $m_{1-c}$ even with unbounded computation, while the sender is also unbounded and learns nothing about $c$, and think about the joint distribution of transcripts.
\end{exercise}

\begin{exercise}
\textbf{Round complexity.} Yao's protocol is \emph{constant round} (the message flow of Figure~\ref{ch34:fig-yao-flow}) while BGW-style secret-sharing MPC uses one communication round per layer of multiplication gates. For a function whose circuit is a single chain of $1000$ multiplications, contrast the round counts of the two approaches, and explain why the Beaver-triple offline phase changes the practical calculus.
\end{exercise}
